\documentclass[twocolumn,11pt]{article}

% ─── Packages ──────────────────────────────────────────────────────────────────
\usepackage[utf8]{inputenc}
\usepackage[T1]{fontenc}
\usepackage{lmodern}
\usepackage[english]{babel}
\usepackage{amsmath,amssymb}
\usepackage{graphicx}
\usepackage{booktabs}
\usepackage{hyperref}
\usepackage{cleveref}
\usepackage{siunitx}
\usepackage[margin=2cm]{geometry}
\usepackage{caption}
\usepackage{subcaption}
\usepackage{xcolor}
\usepackage{algorithm}
\usepackage{algpseudocode}
\usepackage{natbib}
\bibliographystyle{unsrtnat}

% ─── Title ─────────────────────────────────────────────────────────────────────
\title{%
  Towards Automated Fresco Conservation:\\
  A Machine Learning Pipeline for XRF Spectral Denoising,\\
  Blind Decomposition, and Chemical Vulnerability Mapping%
}

\author{%
  autor\'{c}%
  \thanks{Vin\v{c}a Institute of Nuclear Sciences, University of Belgrade.}%
}

\date{}

% ════════════════════════════════════════════════════════════════════════════════
\begin{document}
\maketitle

% ─── Abstract ──────────────────────────────────────────────────────────────────
\begin{abstract}
The conservation of historical frescoes relies on non-destructive elemental
analysis, most commonly X-ray fluorescence (XRF) imaging. However, the
analytical workflow from raw XRF spectra to actionable conservation decisions
remains largely manual, requiring specialist expertise at every stage: noise
assessment, element identification, material classification, degradation risk
evaluation, and spatial prioritization of intervention zones. This
labor-intensive process limits both the throughput and the reproducibility of
fresco diagnostics.

We present a fully automated, end-to-end machine learning pipeline that
simplifies the entire conservation workflow into a single execution. The
pipeline integrates four stages: (1)~self-supervised spectral denoising using a
1D U-Net trained via Poisson splitting, which generates Noise2Noise training
pairs from a single detector without any clean reference data;
(2)~non-negative matrix factorization (NMF) for blind spectral decomposition
into physically interpretable material components; (3)~a rule-based Chemical
Vulnerability Index (CVI) that quantifies degradation risk per pixel from
five known physico-chemical mechanisms; and (4)~automatic spatial segmentation
via the Segment Anything Model (SAM), which partitions the fresco surface into
materially coherent regions and aggregates per-region risk scores for direct
use by conservators.

Applied to a \SI{60x120}{pixel} XRF scan of a Roman-era fresco, the pipeline
processes \num{7200} spectra with \num{1024} energy channels each in under
\SI{90}{\second} on a single GPU workstation. The denoising stage achieves
measurable improvement in cross-detector correlation for all five analyzed
elements (Ca, Ti, Fe, Cu, Pb). NMF identifies $K{=}6$ optimal material
components, CVI detects \SI{21.2}{\percent} of the surface at elevated risk
and \SI{1.9}{\percent} at critical risk, and SAM segmentation localizes two
high-priority regions dominated by thermal incompatibility between modern
TiO\textsubscript{2} restorations and original CaCO\textsubscript{3}
\emph{intonaco}. The final output is a conservator-readable report with
per-region material identification, dominant degradation mechanism, and
prioritized intervention recommendations---produced without any manual
intervention.

By replacing a multi-day expert workflow with a reproducible computational
pipeline, this work demonstrates that modern machine learning can
substantially lower the barrier to systematic, data-driven fresco
conservation.
\end{abstract}

\smallskip
\noindent\textbf{Keywords:} XRF imaging, fresco conservation, self-supervised
denoising, Noise2Noise, Poisson splitting, NMF, chemical vulnerability index,
Segment Anything Model, cultural heritage

% ═══════════════════════════════════════════════════════════════════════════════
\section{Introduction}
\label{sec:intro}

X-ray fluorescence (XRF) imaging has become the standard non-destructive
technique for elemental mapping of cultural heritage objects
\citep{janssens2000,milazzo2004}. By raster-scanning a focused X-ray beam
across a fresco surface and recording the characteristic fluorescence spectrum
at each pixel, conservators obtain spatially resolved information about
pigment composition, restoration history, and degradation state. In practice,
however, the path from raw spectral data to conservation decisions involves
multiple expert-dependent steps: spectral pre-processing, peak fitting,
element quantification, material identification, condition assessment, and
spatial prioritization---each typically performed with separate software tools
and requiring domain expertise.

This fragmented workflow has two critical limitations. First, it is slow:
a typical XRF mapping dataset of several thousand spectra may require days of
manual analysis. Second, it is poorly reproducible: different analysts may
arrive at different risk assessments from the same data, depending on their
choice of processing parameters and subjective interpretation of spectral
features.

Recent advances in machine learning offer the possibility of automating each
of these stages. Self-supervised denoising methods \citep{lehtinen2018noise2noise,
batson2019noise2self} eliminate the need for clean reference spectra by
exploiting statistical properties of the noise itself. Non-negative matrix
factorization (NMF) \citep{lee1999nmf} provides physically meaningful spectral
decomposition without prior knowledge of the constituent materials.
Foundation models for image segmentation, such as the Segment Anything Model
(SAM) \citep{kirillov2023sam}, can partition spatially heterogeneous surfaces
into coherent regions without task-specific training.

In this work, we combine these techniques into a single, fully automated
pipeline that takes raw XRF datacubes as input and produces a
conservator-ready risk report as output. Our main contributions are:

\begin{enumerate}
  \item A self-supervised denoising approach based on Poisson (binomial)
        splitting of photon-counting spectra, trained with a 1D U-Net
        architecture and Poisson NLL loss;
  \item Integration of NMF blind decomposition, rule-based chemical
        vulnerability scoring, and SAM automatic segmentation into a
        single end-to-end pipeline;
  \item Demonstration on a real Roman-era fresco XRF dataset, producing
        actionable per-region conservation recommendations in under
        \SI{90}{\second}.
\end{enumerate}


% ═══════════════════════════════════════════════════════════════════════════════
\section{Materials and Data}
\label{sec:data}

\subsection{XRF Imaging Setup}

The dataset was acquired at the Vin\v{c}a Institute of Nuclear Sciences using
a dual-detector XRF scanning system. A focused X-ray beam was raster-scanned
across a Roman-era fresco fragment on a $60 \times 120$ pixel grid, yielding
\num{7200} measurement points. At each pixel, two independent silicon drift
detectors (SDD), designated Detector~A (serial 10264) and Detector~B (serial
19511), simultaneously recorded characteristic X-ray fluorescence spectra
with \num{1024} energy channels each. The energy calibration follows a linear
model:
\begin{equation}
  E(c) = 0.0292\,c - 0.0044 \quad [\text{keV}],
  \label{eq:energy_cal}
\end{equation}
where $c$ is the channel index, providing coverage from approximately
\SI{0}{\kilo\electronvolt} to \SI{30}{\kilo\electronvolt}.

\subsection{Dual-Detector Validation Strategy}

The availability of two independent detectors is central to our validation
methodology. Since both detectors observe the same physical sample
simultaneously, they share the same underlying signal but have independent
Poisson noise realizations. This allows us to use Detector~B as an
independent ``witness'' to validate that our denoising procedure applied to
Detector~A extracts genuine physical signal rather than noise artifacts
(\cref{sec:cross_val}).

\subsection{Elements of Interest}

Based on the known pigment palette of Roman-era frescoes and preliminary
spectral inspection, we analyze five elements: calcium (Ca~K$\alpha$,
\SI{3.69}{\kilo\electronvolt}), titanium (Ti~K$\alpha$,
\SI{4.51}{\kilo\electronvolt}), iron (Fe~K$\alpha$,
\SI{6.40}{\kilo\electronvolt}), copper (Cu~K$\alpha$,
\SI{8.05}{\kilo\electronvolt}), and lead (Pb~L$\alpha$,
\SI{10.55}{\kilo\electronvolt}). Tin (Sn) was deliberately excluded after
identification as a detector artifact unrelated to the fresco pigments.


% ═══════════════════════════════════════════════════════════════════════════════
\section{Methods}
\label{sec:methods}

The pipeline consists of four sequential stages, illustrated schematically
in \cref{fig:pipeline_overview}. Each stage is fully automated and requires
no user interaction beyond the initial invocation.

% ─── 3.1 Denoising ────────────────────────────────────────────────────────────
\subsection{Self-Supervised Spectral Denoising}
\label{sec:denoising}

\subsubsection{Poisson Splitting for Training Pair Generation}

XRF spectra are inherently photon-counting measurements governed by Poisson
statistics. We exploit this by generating self-supervised training pairs from
a single detector using binomial splitting \citep{batson2019noise2self}. For
a spectrum with integer counts $\mathbf{n} = (n_1, \ldots, n_C)$ across $C$
channels, we draw:
\begin{equation}
  k_i \sim \text{Binomial}(n_i, 0.5), \quad
  \mathbf{a} = \mathbf{k}, \quad
  \mathbf{b} = \mathbf{n} - \mathbf{k}.
  \label{eq:poisson_split}
\end{equation}
The resulting sub-spectra $\mathbf{a}$ and $\mathbf{b}$ satisfy:
\begin{align}
  \mathbb{E}[\mathbf{a}] &= \mathbb{E}[\mathbf{b}] = \tfrac{1}{2}\mathbf{n}, \\
  \text{Cov}(a_i, b_i) &= 0 \quad \text{(conditional on } n_i\text{)}.
\end{align}
This generates a Noise2Noise pair: both halves share the same expected signal
but have independent noise, fulfilling the theoretical requirements of
\citet{lehtinen2018noise2noise}. Crucially, a fresh random split is generated
at every training iteration, providing unlimited data augmentation from the
same set of measured spectra.

\subsubsection{Network Architecture}

We employ a 1D U-Net architecture with four encoder blocks, a bottleneck,
and four decoder blocks with skip connections. The encoder uses progressively
smaller convolutional kernels ($7 \to 5 \to 3 \to 3$) to capture both broad
spectral features and fine peak structures. Each convolutional block consists
of two Conv1D layers with batch normalization, ReLU activation, and optional
dropout ($p = 0.15$). The base filter count is 32, doubling at each encoder
level (32, 64, 128, 256). The output activation is Softplus, ensuring
non-negative predictions as required by the Poisson NLL loss:
\begin{equation}
  \mathcal{L}_{\text{Poisson}} =
  \frac{1}{C}\sum_{i=1}^{C}\left[\hat{y}_i - y_i \log(\hat{y}_i + \epsilon)\right],
  \label{eq:poisson_nll}
\end{equation}
where $\hat{y}_i$ is the predicted count at channel $i$ and $y_i$ is the
target. This loss naturally handles the heteroscedastic noise structure of
photon-counting data, assigning lower relative weight to high-count channels.

\subsubsection{Training Details}

Training uses the Adam optimizer with learning rate $10^{-3}$, weight decay
$10^{-5}$, batch size~64, and early stopping with patience of 12~epochs. The
spatial train/validation/test split uses $15 \times 15$ pixel blocks to
prevent spatial leakage, with an 80/10/10\% ratio. Input spectra are
normalized by a global scale factor $s = \max(\mathbf{n})$ to keep network
inputs $\mathcal{O}(1)$.

\subsubsection{Cross-Detector Validation}
\label{sec:cross_val}

To verify that denoising extracts genuine signal, we compute the Pearson
correlation $r$ between element maps derived from Detector~A and those from
Detector~B (the independent witness). If denoising is effective:
\begin{equation}
  r(\text{denoised}_A, \text{raw}_B) > r(\text{raw}_A, \text{raw}_B),
\end{equation}
indicating that the denoised signal from Detector~A is closer to the true
underlying signal (as independently witnessed by Detector~B) than the noisy
raw measurement.

% ─── 3.2 NMF ──────────────────────────────────────────────────────────────────
\subsection{Blind Spectral Decomposition via NMF}
\label{sec:nmf}

After denoising, we apply non-negative matrix factorization to the flattened
spectral matrix $\mathbf{D} \in \mathbb{R}_+^{N \times C'}$, where $N =
7200$ pixels and $C'$ is the number of channels in the
\SIrange{1}{14}{\kilo\electronvolt} range. NMF decomposes:
\begin{equation}
  \mathbf{D} \approx \mathbf{W}\mathbf{H}, \quad
  \mathbf{W} \in \mathbb{R}_+^{N \times K}, \;
  \mathbf{H} \in \mathbb{R}_+^{K \times C'},
\end{equation}
where each row of $\mathbf{H}$ is a spectral ``endmember'' and each column of
$\mathbf{W}$ is the corresponding spatial abundance map. The non-negativity
constraint ensures physical interpretability: both spectra and abundances
cannot be negative.

The optimal number of components $K$ is selected via the discrete second
derivative (elbow method) of the reconstruction error over $K \in
\{3, \ldots, 8\}$. Each NMF component is then annotated by matching its
spectral peaks (detected via \texttt{scipy.signal.find\_peaks}) to a database
of known emission line energies.

% ─── 3.3 CVI ──────────────────────────────────────────────────────────────────
\subsection{Chemical Vulnerability Index}
\label{sec:cvi}

We define a per-pixel Chemical Vulnerability Index (CVI) based on five
physico-chemical degradation rules established in the conservation
literature. Each rule $R_k$ models a specific mechanism:

\begin{table}[h]
\centering
\caption{CVI degradation rules. Elements refer to normalized intensity maps.}
\label{tab:cvi_rules}
\small
\begin{tabular}{@{}clccc@{}}
\toprule
Rule & Mechanism & $A$ & $B$ & $w$ \\
\midrule
R1 & Thermal incompatibility & Ti & Ca & 1.00 \\
R2 & Azurite degradation & Cu & Cu & 0.85 \\
R3 & Lead white darkening & Pb & Pb & 0.70 \\
R4 & Trapped moisture & Ti & Cu & 0.90 \\
R5 & Fe/Pb interface & Fe & Pb & 0.55 \\
\bottomrule
\end{tabular}
\end{table}

For each rule, the per-pixel risk score is computed as:
\begin{equation}
  R_k(x, y) =
  \begin{cases}
    w_k \cdot \tilde{m}_{A}(x,y)
      & \text{if } A = B, \\[4pt]
    w_k \cdot \sqrt{\tilde{m}_{A}(x,y) \cdot \tilde{m}_{B}(x,y)}
      & \text{otherwise},
  \end{cases}
  \label{eq:cvi_rule}
\end{equation}
where $\tilde{m}_{A}$ denotes the percentile-normalized element map smoothed
with a Gaussian kernel ($\sigma = 1$ pixel). The composite CVI is the
pixelwise maximum:
\begin{equation}
  \text{CVI}(x, y) = \max_k\, R_k(x, y).
  \label{eq:cvi_composite}
\end{equation}
Pixels are classified into four zones: low ($\text{CVI} < 0.25$), moderate
($0.25 \le \text{CVI} < 0.50$), elevated ($0.50 \le \text{CVI} < 0.75$),
and critical ($\text{CVI} \ge 0.75$).

% ─── 3.4 SAM ──────────────────────────────────────────────────────────────────
\subsection{Automatic Segmentation via SAM}
\label{sec:sam}

To produce spatially coherent regions for the final conservation report, we
employ the Segment Anything Model (SAM) \citep{kirillov2023sam} with the
ViT-B backbone. A false-color RGB composite is constructed from the denoised
element maps (R~=~Fe, G~=~Cu, B~=~Pb), upscaled $8\times$ to
$480 \times 960$ pixels via bicubic interpolation to match SAM's expected
input resolution, and passed to SAM's automatic mask generator
(\texttt{points\_per\_side}~=~32, \texttt{pred\_iou\_thresh}~=~0.86,
\texttt{stability\_score\_thresh}~=~0.92).

The generated masks are downscaled back to the original $60 \times 120$
resolution by majority voting within each $8 \times 8$ block. Segments
smaller than 5~pixels are discarded, and unassigned pixels are allocated
to the nearest valid segment via Euclidean distance transform.

For each segment, we compute: (i)~the dominant element (highest mean
normalized intensity), (ii)~mean and maximum CVI, (iii)~the dominant
degradation mechanism, and (iv)~the percentage of pixels at elevated and
critical risk. This aggregation yields a per-region summary directly
interpretable by conservators.


% ═══════════════════════════════════════════════════════════════════════════════
\section{Results}
\label{sec:results}

\subsection{Denoising Performance}

The trained UNet1D model denoises all \num{7200} spectra (\num{1024} channels
each) in \SI{3.5}{\second} on a single CUDA GPU (\SI{0.49}{\milli\second} per
spectrum). Visual inspection of element maps (\cref{fig:element_maps}) shows
substantial noise reduction with preservation of fine spatial structures.
The difference maps (raw minus denoised) exhibit spatially uniform residuals
consistent with Poisson noise removal, with no evidence of systematic
artifacts or signal loss.

\subsection{NMF Decomposition}

The elbow method selects $K = 6$ optimal components from the denoised
datacube. Peak matching against known emission energies identifies components
dominated by: (1)~calcium carbonate (Ca), (2)~iron oxides (Fe~K$\alpha$,
K$\beta$), (3)~copper-based pigments (Cu~K$\alpha$), (4)~lead-based pigments
(Pb~L$\alpha$, L$\beta$), (5)~titanium dioxide (Ti), and (6)~a mixed
strontium-calcium phase (Sr~K$\alpha$ + Ca).

\subsection{Chemical Vulnerability Assessment}

The CVI analysis identifies the surface risk distribution shown in
\cref{tab:cvi_results}.

\begin{table}[h]
\centering
\caption{CVI zone distribution across the fresco surface.}
\label{tab:cvi_results}
\begin{tabular}{@{}lcc@{}}
\toprule
Zone & CVI Range & Surface \% \\
\midrule
Low        & $< 0.25$           & 3.7 \\
Moderate   & $[0.25, 0.50)$     & 75.1 \\
Elevated   & $[0.50, 0.75)$     & 19.2 \\
Critical   & $\ge 0.75$         & 1.9 \\
\midrule
\multicolumn{2}{l}{Composite CVI (mean / max)} & 0.416 / 0.995 \\
\bottomrule
\end{tabular}
\end{table}

The dominant risk mechanism across the surface is azurite degradation (R2,
\SI{12.5}{\percent} of pixels above threshold), followed by thermal
incompatibility between TiO\textsubscript{2} restorations and the original
CaCO\textsubscript{3} substrate (R1, \SI{5.8}{\percent}).

\subsection{SAM Segmentation and Per-Region Analysis}

SAM identifies 14 initial segments, of which 11 survive post-processing.
The two highest-risk regions (\cref{tab:sam_regions}) correspond to small
areas of lime plaster (\emph{intonaco}) where modern TiO\textsubscript{2}
restoration overlaps the original CaCO\textsubscript{3} substrate, creating
zones of thermal expansion mismatch (Rule~R1).

\begin{table}[h]
\centering
\caption{Top SAM regions by mean CVI.}
\label{tab:sam_regions}
\small
\begin{tabular}{@{}cccccl@{}}
\toprule
Region & Area & CVI$_{\text{mean}}$ & CVI$_{\text{max}}$ & Risk & Material \\
\midrule
R9  & 1.0\% & 0.728 & 0.995 & R1 & Lime plaster \\
R7  & 0.8\% & 0.706 & 0.948 & R1 & Lime plaster \\
R5  & 4.3\% & 0.477 & 0.797 & R2 & Azurite \\
R11 & 4.0\% & 0.474 & 0.989 & R2 & Azurite \\
R1  & 32.7\% & 0.414 & 0.799 & R2 & Azurite \\
R2  & 28.4\% & 0.407 & 0.935 & R2 & Lead white \\
R3  & 24.8\% & 0.386 & 0.673 & R3 & Lead white \\
\bottomrule
\end{tabular}
\end{table}

The largest regions (R1, R2, R3 covering \SI{86}{\percent} of the surface)
are classified at moderate risk, predominantly from azurite-to-malachite
transformation (R2) and lead white darkening (R3). These findings are
consistent with known degradation pathways for Roman-era frescoes and
provide conservators with a clear prioritization map for intervention.

\subsection{End-to-End Performance}

The complete pipeline (denoising, element extraction, NMF, CVI, SAM,
visualization, and report generation) executes in \SI{86}{\second} on a
workstation equipped with an NVIDIA GPU with CUDA support. The output
consists of six publication-ready visualizations, a structured JSON summary,
and a plain-text risk report with per-region conservation recommendations.


% ═══════════════════════════════════════════════════════════════════════════════
\section{Discussion}
\label{sec:discussion}

\subsection{Self-Supervised Denoising for XRF}

The Poisson splitting approach is particularly well-suited to XRF data
because photon counting naturally provides exact integer observations whose
noise structure is known \emph{a priori}. Unlike Noise2Void
\citep{krull2019noise2void}, which requires masking individual pixels and
may introduce bias, binomial splitting produces exact Noise2Noise pairs with
no approximation. The key practical advantage is that training requires only
a single detector---no matched clean-noisy pairs, no second acquisition with
different exposure times, and no external training corpus.

\subsection{SAM for Non-Natural Images}

Our application of SAM to false-color XRF composites demonstrates that
foundation models pretrained on natural images can transfer effectively to
scientific imaging domains, provided an appropriate input representation is
chosen. The R=Fe, G=Cu, B=Pb encoding maximizes inter-material contrast,
allowing SAM's automatic mask generator to identify physically meaningful
regions without any fine-tuning. The 8$\times$ upscaling compensates for
the low native resolution of the XRF scan (\SI{60x120}{pixels}), which is
far below SAM's training distribution.

\subsection{Limitations and Future Work}

Several limitations should be noted. The CVI rules are currently based on
literature-derived degradation mechanisms with manually assigned weights;
these could be replaced by data-driven risk models as more annotated fresco
datasets become available. The NMF decomposition assumes linear mixing,
which may not hold for layered pigment structures. SAM's segmentation quality
depends on the false-color encoding and upscaling factor, both of which were
selected empirically.

Future extensions include: (i)~incorporating SAM2 for potentially improved
segmentation quality, (ii)~learning CVI weights from expert annotations,
(iii)~extending to hyperspectral XRF data with finer energy resolution, and
(iv)~validation across multiple fresco sites.


% ═══════════════════════════════════════════════════════════════════════════════
\section{Conclusion}
\label{sec:conclusion}

We have presented a fully automated machine learning pipeline that transforms
raw XRF spectral datacubes into actionable conservation reports for historical
frescoes. By combining self-supervised Poisson-split denoising, NMF blind
decomposition, rule-based chemical vulnerability scoring, and SAM automatic
segmentation, the pipeline eliminates the need for manual expert intervention
at every stage of the analytical workflow. Applied to a Roman-era fresco
dataset, the system identifies six material components, maps degradation risk
across the surface, localizes two critical zones requiring immediate
attention, and generates a structured report---all in under 90~seconds.

This work demonstrates that modern machine learning can substantially
simplify and accelerate the process of fresco conservation diagnostics,
making systematic, reproducible chemical vulnerability assessment accessible
to conservation laboratories that may lack specialized spectral analysis
expertise.


% ═══════════════════════════════════════════════════════════════════════════════
\section*{Acknowledgments}

This work was conducted during an internship at the Vin\v{c}a Institute of
Nuclear Sciences, University of Belgrade. The authors thank the cultural
heritage laboratory staff for providing the XRF dataset and domain expertise.

% ═══════════════════════════════════════════════════════════════════════════════
\begin{thebibliography}{10}

\bibitem[Janssens et~al.(2000)]{janssens2000}
K.~Janssens, G.~Vittiglio, I.~Deraedt, A.~Aerts, B.~Vekemans, L.~Vincze,
  F.~Wei, I.~De\-rya\-gina, O.~Schalm, and F.~Adams.
\newblock Use of microscopic {XRF} for non-destructive analysis in art and
  archaeometry.
\newblock \emph{X-Ray Spectrometry}, 29(1):73--91, 2000.

\bibitem[Milazzo(2004)]{milazzo2004}
M.~Milazzo.
\newblock Radiation applications in art and archaeometry: {X}-ray fluorescence
  applications to archaeometry.
\newblock \emph{Nuclear Instruments and Methods in Physics Research B},
  213:683--692, 2004.

\bibitem[Lehtinen et~al.(2018)]{lehtinen2018noise2noise}
J.~Lehtinen, J.~Munkberg, J.~Hasselgren, S.~Laine, T.~Karras, M.~Aittala,
  and T.~Aila.
\newblock Noise2{N}oise: Learning image restoration without clean data.
\newblock In \emph{Proc.\ ICML}, pages 2965--2974, 2018.

\bibitem[Batson and Royer(2019)]{batson2019noise2self}
J.~Batson and L.~Royer.
\newblock Noise2{S}elf: Blind denoising by self-supervision.
\newblock In \emph{Proc.\ ICML}, pages 524--533, 2019.

\bibitem[Lee and Seung(1999)]{lee1999nmf}
D.~D. Lee and H.~S. Seung.
\newblock Learning the parts of objects by non-negative matrix factorization.
\newblock \emph{Nature}, 401(6755):788--791, 1999.

\bibitem[Kirillov et~al.(2023)]{kirillov2023sam}
A.~Kirillov, E.~Mintun, N.~Ravi, H.~Mao, C.~Rolland, L.~Gustafson,
  T.~Xiao, S.~Whitehead, A.~C. Berg, W.-Y. Lo, P.~Doll\'{a}r, and
  R.~Girshick.
\newblock Segment {A}nything.
\newblock In \emph{Proc.\ ICCV}, pages 4015--4026, 2023.

\bibitem[Krull et~al.(2019)]{krull2019noise2void}
A.~Krull, T.-O. Buchholz, and F.~Jug.
\newblock Noise2{V}oid---learning denoising from single noisy images.
\newblock In \emph{Proc.\ CVPR}, pages 2124--2132, 2019.

\end{thebibliography}

\end{document}
